% GENERATED by md2tex.py from sections/06_related_work.md - edit the .md, not this file
\section{Related Work}

\textbf{Associative capacity of fast-weight states.} \citet{nichani2024understandingfactualrecalltransformers} prove that under argmax decoding a rank-one matrix can
recover on the order of $d$ associations, which is why no result in
this paper derives a rank or capacity claim from argmax recall: the
rank law of Section 3 forces exact continuous recovery through a
Procrustes-aligned cosine readout, and every argmax recall number in
Section 4 carries their caveat explicitly. Their analysis is the reason
the causal razor of Section 3 refuses argmax decoding anywhere a rank
claim depends on it.

\textbf{Associative recall in transformers.} Multi-query associative recall
is a regime where softmax attention is documented as strong: \citet{arora2023zoologymeasuringimprovingrecall} introduce the MQAR benchmark and show attention solves it
where sub-quadratic models struggle; \citet{jelassi2024repeatmetransformersbetter} show
transformers outperform state-space models at copying; \citet{olsson2022incontextlearninginductionheads} show that two-layer transformers reliably develop induction
circuitry from scratch. Our transformer baseline's chance-level reading
therefore runs against the expectation this literature sets, and we do
not interpret it as an architectural inability. The leading candidate
explanation is optimization: the arm trained at the shared default
learning rate with no architecture-specific search on the recall task,
and its training loss is near flat across the run (Section 4.1). The
separation verdict is accordingly carried by the vector-state ablation
comparison alone, with the transformer reading disclosed as a
degenerate-baseline datum; Section 7 states the measurement that would
resolve the question.

\textbf{Rank measurement in linear attention.} Two contemporaneous studies
\citep{nazari2026keystatereductionlinear,sun2026staterankdynamicslinear} measure the rank of
linear-attention states descriptively, on frozen checkpoints. Section 3
differs in both directions of inference: the force-rank arms intervene
at train time, and the necessity/sufficiency step at $d_{\min}$ is a
causal result a descriptive measurement cannot license.

\textbf{Key normalization and stability.} The qk-normalization line (Kimi Linear, \citealp{kimiteam2025kimilinearexpressiveefficient}; Qwen3-Next) conditions
single-vector eigenvalue stability of the state update. The attractor
of Section 5 is a different object, cross-key population geometry. It
was measured with qk normalization active throughout, and removing the
normalization is a within-noise null at $n{=}3$ \% evidence: R7,
so the two phenomena are empirically as well as conceptually distinct.

\textbf{DeltaNet and gated variants.} DeltaNet \citep{yang2025parallelizinglineartransformersdelta} and Gated DeltaNet \citep{yang2025gateddeltanetworksimproving} supply the architecture substrate for the capability and pathology
legs (Sections 4 and 5) and benchmark it for quality and throughput.
This paper asks what the states represent: the rank law of Section 3 is
established on a separate matrix-state encoder family with no
delta-rule write, while the causal recall capability and the
write-induced geometry at scale are established on the delta-rule
substrate itself. The gating
axis reappears here only as a measured factor in Section 5.2, where it
trends toward amplifying the attractor without confirming.

\textbf{Prior negative result in this program.} A published companion study
("The Gradient Does Not See Rank," ICML 2026 MI workshop) showed that
bolting a matrix state onto a pretrained backbone through a
flatten-then-project readout leaves the gradient rank-blind. The
present paper is the from-scratch counterpart: when the state is native
and the readout preserves structure, SGD does see rank (Section 3), and
the resulting medium supports a capability its matched ablations lack
(Section 4).

